Let \(m_O\) denote the number of observational pilot records.  At every center \(P_{\chi,\epsilon_m}(G=O)=1/2\), so
\[
 \frac{m_O}{m}=\frac12+O_P(m^{-1/2}),
 \qquad P(m_O<m/4)\longrightarrow0.                       \tag{53}
\]
Every coordinate product in \(Y^2\), \(\xi_dY\), or \(\xi_d\xi_d^\top\) is bounded uniformly on the binary family.  Expressing an observational average as the ratio of its \(m^{-1}\)-scaled source-weighted sum to \(m_O/m\), Chebyshev's inequality and (53) give, uniformly over the two designs and finitely many coordinates,
\[
 \max_{d\in\mathcal D}\left(
 |\widehat s_{d,m}-s_d|
 +\|\widehat g_{d,m}-g_d\|
 +\|\widehat\Gamma_{d,m}-\Gamma_d\|_{\mathrm{op}}
 \right)=O_P(m^{-1/2}),                                   \tag{54}
\]
where these empirical moments are proof intermediates obtained by expanding the quadratic in Definition~\(\mathrm{def{:}clw\text{-}empirical\text{-}linear\text{-}cost\text{-}selector}\), and the population moments are evaluated at \(P_{\chi,\epsilon_m}\).

Lemma~\(\mathrm{lem{:}sharp\text{-}uniform\text{-}binary\text{-}elbow}\) gives Gram eigenvalues \(1,1\pm xy\).  For \(d_{\mathrm p}\), \(xy=\epsilon_my_\chi<10^{-3}\), and for \(d_{\mathrm e}\), \(xy=1/4\).  Hence
\[
 \min_{d\in\mathcal D}\lambda_{\min}(\Gamma_d)\ge\frac34. \tag{55}
\]
Equations (54)--(55) imply that, with probability tending to one, every empirical Gram is nonsingular with smallest eigenvalue at least \(3/8\).  On this event,
\[
 \widehat r_{d,m}^{\mathrm{lin}}
 =\widehat s_{d,m}-\widehat g_{d,m}^\top
 \widehat\Gamma_{d,m}^{-1}\widehat g_{d,m},
 \qquad
 r_d^{\mathrm{lin}}=s_d-g_d^\top\Gamma_d^{-1}g_d.          \tag{56}
\]
The inverse identity
\[
 \widehat\Gamma_{d,m}^{-1}-\Gamma_d^{-1}
 =\widehat\Gamma_{d,m}^{-1}
 (\Gamma_d-\widehat\Gamma_{d,m})\Gamma_d^{-1}             \tag{57}
\]
and the uniform inverse bounds from (55) show that the difference in (56) is \(O_P(m^{-1/2})\), uniformly in \(d\).  On the exceptional singular event the minimized residual remains in \([0,1]\), since the zero coefficient is feasible and \(Y^2=1\); the exceptional probability tends to zero, so it does not alter the \(O_P(m^{-1/2})\) conclusion.  This proves (1).

Fix \(B\ge5\chi\).  Both floor counts are positive.  Multiplication of (1) by the fixed constants \(C_d/B\) proves the continuous criterion rate, and division by the fixed positive integers \(\lfloor B/C_d\rfloor\) proves (2).

It remains to show that the population gaps do not collapse as \(\epsilon_m\downarrow0\).  The center formulas give
\[
 C_{d_{\mathrm p}}r_{d_{\mathrm p}}^{\mathrm{lin}}(P_{\chi,\epsilon_m})<\frac32,
 \qquad
 C_{d_{\mathrm e}}r_{d_{\mathrm e}}^{\mathrm{lin}}(P_{\chi,\epsilon_m})=3, \tag{58}
\]
so the continuous population gap is bounded below by \(3/(2B)\).  For the integer criterion, set
\[
 k_{\mathrm p}=\left\lfloor\frac{B}{5\chi}\right\rfloor,
 \qquad
 k_{\mathrm e}=\left\lfloor\frac B5\right\rfloor.
\]
Both integers are positive.  If \(z=B/(5\chi)\), then \(z<2\lfloor z\rfloor=2k_{\mathrm p}\), so
\[
 k_{\mathrm e}\le\chi z<2\chi k_{\mathrm p}.
\]
Consequently the fixed number
\[
 \delta_{\chi,B}
 =\frac{3}{5k_{\mathrm e}}-
 \frac{3}{10\chi k_{\mathrm p}}>0.                        \tag{59}
\]
Since
\[
 r_{d_{\mathrm p}}^{\mathrm{lin}}(P_{\chi,\epsilon_m})
 <\frac3{10\chi},
 \qquad
 r_{d_{\mathrm e}}^{\mathrm{lin}}(P_{\chi,\epsilon_m})=\frac35,
\]
the exact integer criterion gap exceeds \(\delta_{\chi,B}\) for every \(m\).  The criterion consistency and the two fixed positive gaps imply
\[
 P_{\chi,\epsilon_m}\!\left\{
 \widehat d_{\mathrm{CLW},m}(B)
 =\widehat d_{\mathrm{CLW},m}^{\mathrm{int}}(B)
 =d_{\mathrm p}\right\}\longrightarrow1.                \tag{60}
\]
At \(B=5\chi\), \(k_{\mathrm p}=1\) and (59) remains strict, so the boundary is included.

On the event in (60), each achieved empirical ratio equals the corresponding population profile ratio at \(d_{\mathrm p}\).  Balanced treatment makes the IKM profile one quarter of the arm-normalized profile design by design, so the ratios agree, and the sharp-elbow lower bound gives
\[
 \frac{\mathcal V_{d_{\mathrm p}}^*(P_{\chi,\epsilon_m})}
 {\mathcal V_{d_{\mathrm e}}^*(P_{\chi,\epsilon_m})}
 =\frac{\mathcal V_{d_{\mathrm p},\mathrm{IKM}}^*(P_{\chi,\epsilon_m})}
 {\mathcal V_{d_{\mathrm e},\mathrm{IKM}}^*(P_{\chi,\epsilon_m})}
 >\frac{\epsilon_m^{-2}+14}{726}\longrightarrow\infty.  \tag{61}
\]
For every finite threshold, (61) eventually exceeds it and the event in (60) has probability tending to one.  This proves (4).  No relation between \(m\) and \(\epsilon_m\) was used.  Each \(\epsilon_m\) is strictly positive, so every law on the path is regular; \(\epsilon=0\) is only a singular limit.  No step uses OGA or a bridge learner.